La versión vulnerable simula un escenario de fallo de red durante la escritura del archivo de salida.
Aunque el sistema no opera sobre una red real, el mecanismo reproduce fielmente los efectos de una
desconexión: datos incompletos, escritura parcial y corrupción del archivo final. El análisis
de esta simulación permite comprender por qué el diseño del proyecto original es resistente ante
este tipo de interrupciones.

\paragraph{Simulación de caída de red.}
En la versión vulnerable, el usuario puede presionar la tecla \texttt{V} durante la ejecución
para activar una bandera \texttt{volatile} llamada \texttt{internetCaido}. Esta bandera se consulta
en múltiples puntos del flujo de escritura, y cuando se activa, el programa abandona el procesamiento
normal e inyecta datos corruptos en el archivo.

\begin{figure}[H]
    \centering
    \begin{minipage}{0.9\textwidth}
    \caption{Hilo de escucha para simular caída de red}
    \label{code:vuln-hilo-red}
    \vspace{-0.2cm}
    \begin{lstlisting}[language=Java, basicstyle=\scriptsize\ttfamily]
private volatile boolean internetCaido = false;

Thread hiloTeclado = new Thread(() -> {
    try {
        while (true) {
            int tecla = System.in.read();
            if (tecla == '\n') { // ENTER: cancelar
                cancelado = true;
                System.exit(0);
            }
            if (tecla == 'v' || tecla == 'V') { // V: simular caida de red
                internetCaido = true;
                System.out.println("\n[!] Internet caido...");
                break;
            }
        }
    } catch (Exception e) {}
});
    \end{lstlisting}
    \end{minipage}
\end{figure}

\paragraph{Escritura en tiempo real sin buffer de integridad.}
El diseño vulnerable escribe los datos al disco \textit{conforme se procesan}, introduciendo
incluso un retardo artificial de 200\,ms entre cada registro para simular una transmisión por red.
Esto significa que si la ``conexión'' se interrumpe a mitad de un alumno, el archivo queda en un
estado parcialmente escrito e irrecuperable.

\begin{figure}[H]
    \centering
    \begin{minipage}{0.9\textwidth}
    \caption{Escritura en tiempo real con comprobaciones de red}
    \label{code:vuln-escritura-red}
    \vspace{-0.2cm}
    \begin{lstlisting}[language=Java, basicstyle=\scriptsize\ttfamily]
private void escribirAlumno(Alumno a, OutputStream out) throws Exception {
    escribirLinea(out, "[" + a.getNombreCompleto() + "]");

    for (DayOfWeek dia : diasSemana()) {
        if (internetCaido) {
            escribirCorrupcionTotal(out);  // 1000 bytes basura
            return;
        }
        escribirLinea(out, "  " + diaEnEspanol(dia) + ":");
        // ...
        for (ClaseHorario c : clases) {
            if (internetCaido) {
                escribirCorrupcionParcial(out);  // 20 bytes basura
                escribirCorrupcionTotal(out);    // 1000 bytes basura
                return;
            }
            escribirLinea(out, "    " + c.getIntervalo() + "  " + c.getNombreMateria());
            Thread.sleep(200);  // Simula latencia de red
        }
    }
}
    \end{lstlisting}
    \end{minipage}
\end{figure}

La tabla~\ref{tab:vuln-red-escenarios} resume los posibles estados del archivo dependiendo del
momento en que se produce la ``caída de red''.

\begin{table}[H]
\centering
\caption{Estados del archivo según el momento de la interrupción de red}
\label{tab:vuln-red-escenarios}
\begin{tabular}{|p{3.8cm}|p{4.5cm}|p{4.8cm}|}
\hline
\textbf{Momento de la caída} & \textbf{Contenido del archivo} & \textbf{Consecuencia} \\
\hline
Antes de procesar algún alumno &
Solo encabezado + ``DATOS INCOMPLETOS'' + 1000 bytes binarios &
Archivo ilegible \\
\hline
A mitad de un día de un alumno &
Alumno parcial + 1000 bytes binarios &
Datos incompletos, archivo corrupto \\
\hline
Entre un registro y otro dentro de un día &
Registro parcial + 20 bytes binarios + 1000 bytes binarios &
Mezcla de texto y basura binaria \\
\hline
\end{tabular}
\end{table}

\paragraph{Protección en el proyecto original.}
El proyecto original no presenta esta vulnerabilidad por dos razones de diseño:

\begin{enumerate}
    \item \textbf{Separación entre procesamiento y escritura.}
    En el proyecto original, primero se cargan \textit{todos} los alumnos en memoria mediante
    \texttt{cargarAlumnosDesdePDF()}, y solo después de verificar que la carga fue exitosa se
    procede a escribir el archivo completo con \texttt{escribirHorariosAlumnos()}. Si cualquier
    PDF falla durante la carga, el proceso se aborta antes de haber escrito un solo byte en
    el archivo de salida.

    \item \textbf{Escritura atómica mediante \texttt{BufferedWriter}.}
    Al utilizar un \texttt{BufferedWriter} envuelto en un bloque \texttt{try-with-resources},
    los datos se escriben de forma controlada y el archivo se cierra correctamente incluso si
    ocurre una excepción. No existe un punto intermedio donde el flujo pueda quedar a medias.
\end{enumerate}

\begin{figure}[H]
    \centering
    \begin{minipage}{0.9\textwidth}
    \caption{Flujo seguro del proyecto original: primero carga, después escribe}
    \label{code:secure-flujo-red}
    \vspace{-0.2cm}
    \begin{lstlisting}[language=Java, basicstyle=\scriptsize\ttfamily]
// Proyecto original: carga completa antes de escribir
List<Alumno> alumnos;
try {
    alumnos = cargarAlumnosDesdePDF();  // Si falla, se aborta aqui
} catch (RuntimeException e) {
    System.err.println("[!] Proceso terminado.");
    return;  // No se escribe nada en disco
}

// Solo si la carga fue exitosa, se escribe el archivo completo
try (BufferedWriter writer = Files.newBufferedWriter(
        rutaSalida,
        StandardOpenOption.CREATE,
        StandardOpenOption.TRUNCATE_EXISTING)) {
    escribirHorariosAlumnos(alumnos, writer);
    writer.newLine();
    escribirHorariosLibres(alumnos, writer);
}
    \end{lstlisting}
    \end{minipage}
\end{figure}

En resumen, la arquitectura del proyecto original garantiza que el archivo de salida
solo se genera cuando \textit{todos} los datos han sido validados y procesados correctamente.
No existe un estado intermedio en el que el archivo pueda quedar parcialmente escrito ni
un mecanismo externo que pueda inyectar contenido corrupto durante la escritura.

\paragraph{Autenticación básica de acceso.}
Como medida de control de acceso, el sistema implementa un mecanismo de autenticación por
usuario y contraseña antes de iniciar cualquier procesamiento. Esto impide que usuarios no
autorizados ejecuten el programa y accedan a la información de los horarios.

La autenticación se realiza mediante el método \texttt{autenticar()} en la clase
\texttt{Generador}, que es lo primero que se invoca al ejecutar el sistema:

\begin{figure}[H]
    \centering
    \begin{minipage}{0.9\textwidth}
    \caption{Método de autenticación con usuario y contraseña}
    \label{code:secure-auth}
    \vspace{-0.2cm}
    \begin{lstlisting}[language=Java, basicstyle=\scriptsize\ttfamily]
private boolean autenticar() {
    Console consola = System.console();
    if (consola == null) {
        System.err.println("[!] No se puede acceder a la consola.");
        return false;
    }

    String usuario = consola.readLine("Usuario: ");
    if (usuario == null || !usuario.equals(USUARIO_CORRECTO)) {
        System.err.println("[!] Usuario incorrecto. Acceso denegado.");
        return false;
    }

    char[] password = consola.readPassword("Contrasena: ");
    if (password == null || !new String(password).equals(PASSWORD_CORRECTA)) {
        System.err.println("[!] Contrasena incorrecta. Acceso denegado.");
        Arrays.fill(password, '0');
        return false;
    }

    Arrays.fill(password, '0');
    System.out.println("[DEBUG] Autenticacion exitosa. Bienvenido, " + usuario);
    return true;
}
    \end{lstlisting}
    \end{minipage}
\end{figure}

El diseño del método incorpora tres decisiones de seguridad relevantes. Primero, se utiliza
\texttt{Console.readPassword()} para leer la contraseña, lo que evita que esta se muestre
en pantalla mientras el usuario la escribe ---a diferencia de un \texttt{Scanner} estándar,
que sí la mostraría. Segundo, si el usuario es incorrecto, el sistema deniega el acceso
inmediatamente sin llegar a solicitar la contraseña, reduciendo la superficie de ataque.
Tercero, una vez verificada, la contraseña se sobreescribe en memoria con
\texttt{Arrays.fill(password, '0')}, eliminando el valor sensible del heap antes de que
el recolector de basura pueda exponerlo.

\begin{table}[H]
\centering
\caption{Comparación de mecanismos de lectura de contraseña}
\label{tab:secure-auth-comparacion}
\begin{tabular}{|p{4cm}|p{4.5cm}|p{4cm}|}
\hline
\textbf{Aspecto} & \textbf{\texttt{Console.readPassword()}} & \textbf{\texttt{Scanner}} \\
\hline
Visibilidad en pantalla & No muestra caracteres & Muestra la contraseña \\
\hline
Tipo de retorno & \texttt{char[]} (borrable) & \texttt{String} (inmutable) \\
\hline
Limpieza de memoria & Posible con \texttt{Arrays.fill} & Imposible \\
\hline
Adecuado para producción & Sí & No \\
\hline
\end{tabular}
\end{table}

El flujo completo de \texttt{ejecutar()} queda así protegido: si la autenticación falla,
el método retorna inmediatamente y ningún PDF es procesado ni ningún archivo es escrito,
garantizando que el acceso a los datos de los alumnos esté siempre condicionado a una
identidad verificada.
